\section{Implementation profiles and claim traceability}\label{sec:implementation}

\subsection{Coordinatewise CLI}
The installable Python package exposes inspection, normalization, target verification, and certificate checking. Its runtime needs only the standard library. Verification constructs a source expression and produces a rewrite trace. Checking reparses the source and replays the supplied trace without invoking normalization search. The two paths share primitive-rule code.

The certificate format remains \texttt{qkf-rewrite-v3}, API results use \texttt{qkf-result-v1}, and the semantic contract is
\begin{quote}\small\ttfamily\raggedright
total-transfer-words-v2; positive-width;\\
modular-arithmetic; SMT-total-div-shift;\\
standard-counts; nonbottom-KB.
\end{quote}
The accepted final coordinatewise targets remain AND, OR, and XOR. A resource limit or unsupported residual produces an explicit unresolved outcome. SHA-256 fingerprints bind the complete supplied source-label map, including helper definitions; they identify bytes, not authorship.

The original three Nice to Meet You examples retain their upstream provenance and MIT notices. AND normalizes to $(z_a\mathbin{|}z_b,o_a\mathbin{\&}o_b)$; OR has the dual optimal table. XOR expands the actual meet helper and uses a checked disjoint-support addition to recover its optimal table. The checker does not assign semantics from the name of a helper.

\subsection{Research execution and summary profiles}
The distribution preserves the earlier \texttt{research/knownbits} and \texttt{research/graal} capsules and their source identities. Replay restores saved evidence and checks it in fresh processes; regenerating search, native Java correspondence, or measurements is a separate action with separately documented dependencies. The historical POSIX capsules retain their environment requirements. The newer source-query and summary modules are tested on Python 3.10 and 3.12 and require a full checkout; they add no runtime dependency to the coordinatewise package.

The principal new entry points are source-interface inference and execution (Section~\ref{sec:sourceinterfaces}), exact pure-row and typed ground queries, guarded source consumers, and \texttt{research.applicable\_summary} (Section~\ref{sec:consumerproofs}). The SDK retains signed \texttt{no\_lemmas} and local-phase \texttt{direct\_cell} defaults; derived-lemma reuse and forcing remain explicit options. The controlled comparisons did not silently change these defaults.

Certificate kinds have distinct scopes. \texttt{qkf-ground-query-certificate-v1} proves statements from a declared finite ground presentation. \texttt{qkf-source-lemma-packet-v1} additionally binds native premises and checked imports to a source context. The summary root \texttt{qkf-applicable-summary-v1}, transported by \texttt{qkf-json-proof-dag-v1}, supplies the evidence used by an applicable SDK object. Preparation and direct-emission changes preserve these public formats and the independent public checker. A valid certificate of one kind cannot be substituted for an obligation belonging to another.

Positive fresh replay excludes producer/search/SMT imports and does not execute the complete source program; native premises are still independently checked. A concrete refutation separately evaluates its supplied witness. Model size, source-state limits, proof-DAG limits, and runtime-width limits are implementation budgets with explicit failure outcomes, not premises of an extrapolation from tested widths. Immutable checked objects protect ordinary API use, not a hostile Python process capable of reflection or monkeypatching.

\subsection{Formal projects and stable claims}
The original \texttt{research/lean}, later \texttt{research/lean\_targets}, and new \texttt{formal/ground} and \texttt{formal/composition} projects pin Lean separately from Python. F1b has a local dependency on F1a. Their build/audit scripts, exporters, positive examples, and negative controls preserve the scope described in Section~\ref{sec:lean}. Table~\ref{tab:traceability} retains the original claim identifiers and identifies the additions separately.

\begin{table}[htbp]\small
\begin{tabularx}{\textwidth}{@{}>{\raggedright\arraybackslash}p{.12\textwidth}>{\raggedright\arraybackslash}X>{\raggedright\arraybackslash}X@{}}
\toprule
Claim & Artifact & Scope\\
\midrule
K1--K2 & Paper supplement: finite core and relation-matrix oracle & Exact finite ground interface and independent finite checks.\\
K3 & Sparse prototype and representation theorem & Conditional coverage/fixed-point contract; historical pass-limit caveat retained.\\
K4--K5 & \texttt{src/qkf\_certifier} & Coordinatewise source replay and nine-row semantics.\\
K6--K7 & Observation/row kernels and original Lean gluing & Labeled consumer observations and arbitrary-length composition.\\
K8--K9 & Retained upper/lower Graal capsules & Exact upper model and conditional lower preservation; separate source contracts.\\
K10 & \texttt{research/lean/QKF} & Original ascending mathematical fragment and nonsign cell model.\\
New source profiles & \texttt{research/inference}, signed interfaces, \texttt{create\_joint} & Restricted source-derived factors and caller contracts with explicit Python trust boundaries.\\
New proof profiles & Pure rows, ground queries, source lemmas, summary SDK & Exact finite algebraic evidence and scoped source consumers.\\
F1a/F1b & \texttt{formal/ground}, \texttt{formal/composition} & Decoded typed positive calculus and conservative checked-lemma composition.\\
\bottomrule
\end{tabularx}
\caption{Stable claims retain their original meaning; new profiles state their additional scope explicitly.}\label{tab:traceability}
\end{table}

Artifact links identify the technical snapshot; measurement reports identify their individual engines and environments. Historical native executions, fresh certificate replay, formal builds, and release packaging checks are separate records. Bundling an experiment does not promote it to a larger source or semantic scope.
